
% Abstract page describing the project scope, methodology, and key findings

\chapter*{ABSTRACT}
\phantomsection
\addcontentsline{toc}{chapter}{ABSTRACT}

Thapathali is one of Kathmandu's busiest intersections, and its signal operation still depends heavily on fixed cycles and manual judgement. That approach is manageable when demand is predictable, but it becomes weak when one approach is overloaded while another has already cleared. This project studies whether a Deep Q-Network (DQN) controller can use the observed traffic state to make better phase decisions for this specific junction.

The simulation network was prepared in SUMO using OpenStreetMap geometry and traffic demand derived from five days of Origin-Destination (OD) and vehicle-count data for Thapathali. At every 15-second decision step, the agent observes a 21-dimensional state vector consisting of PCU-weighted loads on incoming and outgoing lanes, maximum stopline waiting times, and the active signal phase. The action space was kept deliberately small: the controller either keeps the current phase (STAY) or moves to the next phase (NEXT) in a three-phase cycle, with a yellow transition before switching. A Max-Pressure based reward was used so that the agent is encouraged to clear loaded approaches without ignoring accumulated waiting time. The agent was trained for 300 episodes with experience replay and epsilon-greedy exploration. When evaluated against a fixed 60-second signal plan on Day 5 demand, the DQN policy reduced average waiting time by 22.90\%, reduced average queue length by 30.06\%, and improved throughput efficiency by 32.65\%. The strongest gains appeared during low-to-moderate demand periods, where a fixed cycle wastes green time. During peak saturation the improvement was smaller, which suggests that signal optimisation helps but cannot remove the capacity limits caused by intersection geometry and heavy demand. Overall, the work provides a simulation-based basis for further study of adaptive signal control at Thapathali, especially if future versions include pedestrians, public transport, and neighbouring intersections.
\vspace{2em}\\
\noindent\textbf{Keywords:} \textit{Adaptive Control, \gls{dqn}, \gls{drl}, Max-Pressure Theory, Queue Length, \gls{sumo}, Throughput, Traffic Signal Control}
